
\documentclass{article}
\usepackage{colortbl}
\usepackage{makecell}
\usepackage{multirow}
\usepackage{supertabular}

\begin{document}

\newcounter{utterance}

\centering \large Interaction Transcript for game `cladder', experiment `full\_v1.5\_default', episode 3752 with selfrefine{-}Qwen3.5{-}9B.
\vspace{24pt}

{ \footnotesize  \setcounter{utterance}{1}
\setlength{\tabcolsep}{0pt}
\begin{supertabular}{c@{$\;$}|p{.15\linewidth}@{}p{.15\linewidth}p{.15\linewidth}p{.15\linewidth}p{.15\linewidth}p{.15\linewidth}}
    \# & \multicolumn{2}{c}{Player} && \multicolumn{2}{c}{Game Master} \\
    \hline

    \theutterance \stepcounter{utterance}  
    & & & \multicolumn{4}{p{0.6\linewidth}}{
        \cellcolor[rgb]{0.9,0.9,0.9}{
            \makecell[{{p{\linewidth}}}]{
                \texttt{\tiny{[P1$\langle$GM]}}
                \texttt{You are an expert in causal inference. The following question is not a typical commonsense query, but rather a meticulously designed question created by a professor specializing in causal inference, intended to assess the students' mastery of the course content.} \\
\texttt{Imagine a self{-}contained, hypothetical world with only the following conditions, and without any unmentioned factors or causal relationships: Listening to jazz has a direct effect on elite institution admission status. Effort has a direct effect on elite institution admission status. For students accepted to elite institutions, the correlation between listening to jazz and being hard{-}working is {-}0.09. If we look at students accepted to elite institutions, does it mean that listening to jazz affects being hard{-}working?} \\
\texttt{Start your answer with 'Yes' or 'No', followed by additional reasoning or evidence to support your explanation.} \\
            }
        }
    }
     \\ \\

    \theutterance \stepcounter{utterance}  
    & & \multicolumn{4}{p{0.6\linewidth}}{
        \cellcolor[rgb]{0.95,0.95,0.95}{
            \makecell[{{p{\linewidth}}}]{
                \texttt{\tiny{[P1$\leftrightarrow$P]}}
                \texttt{{[}DRAFT{]}} \\
\texttt{No} \\
\\ 
\texttt{**Reasoning:**} \\
\\ 
\texttt{In this hypothetical world, we are dealing with a classic case of **selection bias** (also known as collider stratification bias). Let us define the variables:} \\
\texttt{*   \$J\$: Listening to jazz} \\
\texttt{*   \$E\$: Effort (being hard{-}working)} \\
\texttt{*   \$A\$: Admission to an elite institution} \\
\\ 
\texttt{The causal structure is defined as:} \\
\texttt{1.  \$J \textbackslash{}rightarrow A\$ (Jazz affects Admission)} \\
\texttt{2.  \$E \textbackslash{}rightarrow A\$ (Effort affects Admission)} \\
\texttt{3.  There is no direct arrow between \$J\$ and \$E\$.} \\
\\ 
\texttt{The variable \$A\$ (Admission) acts as a **collider** because it is a common effect of both \$J\$ and \$E\$. In causal inference, conditioning on a collider (i.e., restricting our analysis to only those students who were accepted, \$A=1\$) induces a spurious statistical association between its parents (\$J\$ and \$E\$), even if they are causally independent in the general population.} \\
\\ 
\texttt{The problem states that among accepted students, the correlation between jazz and effort is {-}0.09. This negative correlation is an artifact of the selection process:} \\
\texttt{*   Students who listen to jazz (\$J=1\$) are less likely to be admitted unless they compensate with high effort (\$E=1\$), because \$J\$ has a direct effect on admission (implied to be negative or simply distinct from effort, creating a trade{-}off in the selection filter).} \\
\texttt{*   Conversely, students who do not listen to jazz (\$J=0\$) need less effort to achieve admission.} \\
\texttt{*   When we look *only* at the accepted group, we are observing a subset where high \$J\$ is paired with high \$E\$ to satisfy the admission threshold, creating a negative correlation between the two traits within this specific subgroup.} \\
\\ 
\texttt{However, the question asks if listening to jazz **affects** being hard{-}working. Causality requires a direct mechanism or a path from \$J\$ to \$E\$ that is not blocked by conditioning. In this world, there is no arrow from \$J\$ to \$E\$. The observed correlation is purely non{-}causal, resulting from the fact that we are looking at a selected sample (the collider). Therefore, the correlation does not imply that listening to jazz causes students to be harder working (or less hard working); it merely reflects the mathematical constraints of the selection criteria.} \\
            }
        }
    }
    & \\ \\

    \theutterance \stepcounter{utterance}  
    & & \multicolumn{4}{p{0.6\linewidth}}{
        \cellcolor[rgb]{0.95,0.95,0.95}{
            \makecell[{{p{\linewidth}}}]{
                \texttt{\tiny{[P1$\leftrightarrow$P]}}
                \texttt{{[}SELF{-}FEEDBACK{]}} \\
\texttt{FORMAT: ok} \\
\texttt{RULES: ok} \\
\texttt{STRATEGY: ok} \\
\texttt{REASON: The answer correctly identifies the causal structure, explains the collider bias mechanism, and logically concludes that the observed correlation does not imply causation.} \\
            }
        }
    }
    & \\ \\

    \theutterance \stepcounter{utterance}  
    & \multicolumn{4}{p{0.6\linewidth}}{
        \cellcolor[rgb]{0.561,0.737,0.561}{
            \makecell[{{p{\linewidth}}}]{
                \texttt{\tiny{[P1$\rangle$GM]}}
                \texttt{No} \\
\\ 
\texttt{**Reasoning:**} \\
\\ 
\texttt{In this hypothetical world, we are dealing with a classic case of **selection bias** (also known as collider stratification bias). Let us define the variables:} \\
\texttt{*   \$J\$: Listening to jazz} \\
\texttt{*   \$E\$: Effort (being hard{-}working)} \\
\texttt{*   \$A\$: Admission to an elite institution} \\
\\ 
\texttt{The causal structure is defined as:} \\
\texttt{1.  \$J \textbackslash{}rightarrow A\$ (Jazz affects Admission)} \\
\texttt{2.  \$E \textbackslash{}rightarrow A\$ (Effort affects Admission)} \\
\texttt{3.  There is no direct arrow between \$J\$ and \$E\$.} \\
\\ 
\texttt{The variable \$A\$ (Admission) acts as a **collider** because it is a common effect of both \$J\$ and \$E\$. In causal inference, conditioning on a collider (i.e., restricting our analysis to only those students who were accepted, \$A=1\$) induces a spurious statistical association between its parents (\$J\$ and \$E\$), even if they are causally independent in the general population.} \\
\\ 
\texttt{The problem states that among accepted students, the correlation between jazz and effort is {-}0.09. This negative correlation is an artifact of the selection process:} \\
\texttt{*   Students who listen to jazz (\$J=1\$) are less likely to be admitted unless they compensate with high effort (\$E=1\$), because \$J\$ has a direct effect on admission (implied to be negative or simply distinct from effort, creating a trade{-}off in the selection filter).} \\
\texttt{*   Conversely, students who do not listen to jazz (\$J=0\$) need less effort to achieve admission.} \\
\texttt{*   When we look *only* at the accepted group, we are observing a subset where high \$J\$ is paired with high \$E\$ to satisfy the admission threshold, creating a negative correlation between the two traits within this specific subgroup.} \\
\\ 
\texttt{However, the question asks if listening to jazz **affects** being hard{-}working. Causality requires a direct mechanism or a path from \$J\$ to \$E\$ that is not blocked by conditioning. In this world, there is no arrow from \$J\$ to \$E\$. The observed correlation is purely non{-}causal, resulting from the fact that we are looking at a selected sample (the collider). Therefore, the correlation does not imply that listening to jazz causes students to be harder working (or less hard working); it merely reflects the mathematical constraints of the selection criteria.} \\
            }
        }
    }
    & & \\ \\

    \theutterance \stepcounter{utterance}  
    & & & \multicolumn{2}{p{0.3\linewidth}}{
        \cellcolor[rgb]{0.9,0.9,0.9}{
            \makecell[{{p{\linewidth}}}]{
                \texttt{\tiny{[GM$|$GM]}}
                \texttt{No} \\
            }
        }
    }
    & & \\ \\

    \theutterance \stepcounter{utterance}  
    & & & \multicolumn{2}{p{0.3\linewidth}}{
        \cellcolor[rgb]{0.9,0.9,0.9}{
            \makecell[{{p{\linewidth}}}]{
                \texttt{\tiny{[GM$|$GM]}}
                \texttt{no} \\
            }
        }
    }
    & & \\ \\

    \theutterance \stepcounter{utterance}  
    & & & \multicolumn{2}{p{0.3\linewidth}}{
        \cellcolor[rgb]{0.9,0.9,0.9}{
            \makecell[{{p{\linewidth}}}]{
                \texttt{\tiny{[GM$|$GM]}}
                \texttt{game\_result = WIN} \\
            }
        }
    }
    & & \\ \\

\end{supertabular}
}

\end{document}
